\section{Anexos}
\label{sec:anexos}

Esta sección reúne el material de apoyo exigido por la guía de la Entrega Final que no encaja como
prosa dentro de los capítulos anteriores. Donde el material solicitado no existe todavía de forma real,
se declara la brecha explícitamente en vez de fabricarlo --- la misma política de integridad aplicada en
el resto de este informe (Sección~\ref{sec:discusion}).

\subsection{Anexo A --- Cadena de búsqueda de la revisión de trabajos relacionados}
\label{anexo:busqueda}

\textbf{Brecha declarada.} La búsqueda dirigida de la Sección~\ref{sec:trabajos-relacionados} reporta el
número de consultas ejecutadas (24) y el número de resultados aproximado (\textasciitilde200), pero
\textbf{no se registraron las cadenas de búsqueda booleanas exactas ni la fecha de cada consulta} en su
momento --- así lo declara \texttt{docs/checklists/PRISMA-2020.md} en el criterio ``Estrategia de
búsqueda completa y reproducible'': \textit{``se reporta el número de consultas (24) pero no las
cadenas de búsqueda exactas ni las fechas de ejecución --- no es reproducible palabra por palabra''}. No
se reconstruye aquí una cadena de búsqueda aproximada a partir de memoria, porque presentarla como si
fuera la consulta real ejecutada sería exactamente el tipo de dato fabricado que este informe corrige en
otros lugares (Secciones~\ref{sec:evaluacion} y \ref{sec:discusion}). Quedan como palabras clave
efectivamente usadas, sin cadena booleana ni fecha verificable: \textit{software requirements
traceability}, \textit{JWT authentication Spring Boot}, \textit{stored procedures vs ORM performance},
\textit{System Usability Scale}, \textit{OWASP dependency scanning}, \textit{CI/CD reproducibility},
entre otras relacionadas con cada uno de los 13 trabajos comparados en la Tabla de la
Sección~\ref{sec:trabajos-relacionados}.

\subsection{Anexo B --- Protocolo del instrumento SUS (System Usability Scale)}
\label{anexo:sus}

El instrumento sigue la forma estándar de Brooke~\cite{brooke1996}: diez enunciados de alternancia de
polaridad (impares en sentido positivo, pares en sentido negativo), cada uno respondido en una escala
Likert de 1 (totalmente en desacuerdo) a 5 (totalmente de acuerdo).

\begin{enumerate}[nosep]
    \item Creo que me gustaría utilizar este sistema con frecuencia.
    \item Encontré el sistema innecesariamente complejo.
    \item Pensé que el sistema era fácil de usar.
    \item Creo que necesitaría el apoyo de un técnico para poder utilizar este sistema.
    \item Encontré que las diversas funciones del sistema estaban bien integradas.
    \item Pensé que había demasiada inconsistencia en el sistema.
    \item Imagino que la mayoría de las personas aprenderían a usar el sistema muy rápidamente.
    \item Encontré el sistema muy pesado/incómodo de usar.
    \item Me sentí muy confiado/seguro al usar el sistema.
    \item Necesité aprender muchas cosas antes de poder empezar a usar el sistema.
\end{enumerate}

\textbf{Cálculo de la puntuación:} para los ítems impares (1,3,5,7,9) se resta 1 a la respuesta; para los
pares (2,4,6,8,10) se resta la respuesta de 5. La suma de las diez contribuciones se multiplica por 2.5,
dando un puntaje individual entre 0 y 100 (no es un porcentaje). El script
\texttt{scripts/sus-analysis.ipynb} implementa exactamente esta fórmula, trae tres casos de autotest que
pasan, y calcula el resultado del grupo directamente desde
\texttt{docs/mediciones/sus/sus-respuestas.csv} (ejecutado de punta a punta con \texttt{nbconvert}, sin
errores).

\textbf{Estado real (Sección~\ref{sec:evaluacion}):} el instrumento se aplicó a 15 personas, pero
\textbf{solo 4 tienen fecha de aplicación verificable} (\textbf{$n=4$}); las otras 11 tienen una fecha
escrita a mano que no se sostiene (posterior al commit que las versiona y, en varios casos, posterior al
día de hoy) y quedan declaradas pendientes de confirmar, no fabricadas ni descartadas --- ver
\texttt{docs/mediciones/sus/SUS-RESULTS.md} para el detalle completo del hallazgo. El protocolo de
consentimiento informado individual está redactado como plantilla en
\texttt{docs/etica/consentimientos/plantilla-consentimiento.md} (explica qué se recoge, para qué, cuánto
tiempo se conserva, quién lo verá y cómo se retira el consentimiento), pero lo que se recolectó en la
práctica fue una nota impresa de consentimiento implícito en cada hoja, no una firma ni una aprobación
individual fechada --- brecha declarada, no subsanada retroactivamente. Una versión anterior de este
proyecto publicó además un puntaje de usabilidad fabricado con diez evaluadores inexistentes; esa cifra
fue retirada explícitamente y se documenta en \texttt{docs/mediciones/sus/SUS-RESULTS.md} y en la
Sección~\ref{sec:evaluacion} de este informe.

\subsection{Anexo C --- Corridas de integración continua}
\label{anexo:ci}

La guía pide capturas de tres corridas verdes de CI. Se optó por una forma de evidencia
\textbf{más fuerte que una imagen}: los datos crudos de la API pública de GitHub, que cualquier tercero
puede volver a consultar con el mismo comando y obtener el mismo resultado, mientras que una captura de
pantalla sólo se puede creer o no creer. Los JSON de respuesta están versionados sin editar en
\texttt{docs/evidencias/ci/}, junto con un resumen legible en su \texttt{README.md}.

Comando de reproducción (no requiere autenticación, el repositorio es público):

\begin{lstlisting}[language=bash,style=codigobase,caption={Consulta de las corridas de CI contra la API publica de GitHub.},label={lst:ci-api}]
REPO=carla22072004/PFC-Presustentaciones-2026
curl -s "https://api.github.com/repos/$REPO/actions/runs?per_page=8"
curl -s "https://api.github.com/repos/$REPO/actions/runs/<RUN_ID>/jobs"
\end{lstlisting}

Las tres corridas capturadas el 2026-09-05 (Tabla~\ref{tab:ci-runs}) corresponden a tres commits
distintos de la rama \texttt{main} y las tres terminaron en \texttt{success}, con sus \textbf{tres jobs}
y la totalidad de sus pasos en verde --- 35 pasos por corrida, ninguno fallido ni omitido.

\begin{table}[H]
\centering
\small
\begin{tabular}{p{2.6cm}p{1.7cm}p{2.2cm}p{6.4cm}}
\toprule
\textbf{Run ID} & \textbf{Commit} & \textbf{Conclusión} & \textbf{Commit defendido} \\
\midrule
33978578357 & \texttt{99a6636} & \texttt{success} & Cobertura de cierre (63.17\,\%, 395 tests) \\
33949588592 & \texttt{581d45d} & \texttt{success} & Diagramas C4, listados y anexos \\
33944802465 & \texttt{76250b7} & \texttt{success} & Merge de la rama del equipo \\
\bottomrule
\end{tabular}
\caption{Corridas verdes de integración continua verificadas contra la API de GitHub.}
\label{tab:ci-runs}
\end{table}

Los tres jobs que ejecuta cada corrida son sustantivos, no decorativos: \textbf{Backend} levanta
Postgres y Redis reales como \textit{service containers} (no mocks), compila, corre la suite completa
con JaCoCo y publica cobertura y resultados de surefire como artefactos, y después corre SpotBugs +
find-sec-bugs; \textbf{Frontend} hace el build de producción de Angular; y \textbf{Entorno} regenera
\texttt{docs/entorno/versions.txt} con las versiones reales de la máquina de CI, para que las versiones
declaradas en este informe no dependan de la máquina de un integrante. Que las tres estén en verde
significa que, en una máquina limpia y contra una base de datos real, el proyecto compila, sus 395
pruebas pasan y el análisis estático de seguridad corre sin romper el pipeline.

\subsection{Anexo D --- Checklist de estándares empíricos de Ralph et al. (2021)}
\label{anexo:ralph}

Además de \texttt{docs/checklists/RALPH-2021-GENERAL.md} (Estándar General, resumido en la
Tabla~\ref{tab:ralph-general}), el repositorio aplica dos checklists adicionales de la misma familia
--- \texttt{RALPH-2021-BENCHMARKING.md} (para las mediciones de rendimiento con k6) y
\texttt{RALPH-2021-ENGINEERING-RESEARCH.md} (para el proceso de desarrollo como tal) --- referenciados
desde la Sección~\ref{sec:marco-teorico} pero no reproducidos aquí en extenso por brevedad.

\begin{table}[H]
\centering
\small
\begin{tabular}{p{7.5cm}p{1.3cm}p{5cm}}
\toprule
\textbf{Criterio (Estándar General)} & \textbf{Cumple} & \textbf{Evidencia resumida} \\
\midrule
Afirmaciones empíricas basadas en datos, no en opinión & Sí & Toda cifra citada tiene archivo crudo de origen (\texttt{DATA-PROVENANCE.md}) \\
Método de recolección reproducible (comandos documentados) & Sí & Comandos exactos en \texttt{k6/README.md}, \texttt{LIGHTHOUSE-REPORT.md}, \texttt{OWASP-AUDIT.md} \\
Contexto/entorno de medición declarado & Sí & \texttt{docs/entorno/versions.txt} + notas metodológicas explícitas \\
Se reportan resultados negativos, no solo favorables & Sí & Endpoint inexistente en k6, 233 hallazgos find-sec-bugs, 14 vulnerabilidades npm sin corregir \\
Correlación vs. causalidad distinguidas & Parcial & Cache fría/caliente sí aísla variable; runs 1--2 de k6 no lo hacen, y el reporte lo aclara \\
Amenazas a la validez discutidas sistemáticamente & No & Sin sección dedicada por reporte individual (sí existe una transversal en Sección~\ref{sec:amenazas-validez} de este informe) \\
Tamaño de muestra justificado & Parcial & $n=30$ en caché sigue convención CLT sin justificación explícita en texto; 5 corridas k6 siguen el mínimo de la guía, no un cálculo de poder \\
Artefactos disponibles para verificación independiente & Sí & \texttt{DATA-PROVENANCE.md} enlaza cada archivo crudo referenciado \\
\bottomrule
\end{tabular}
\caption{Checklist de estándares empíricos generales de Ralph et al.~\cite{ralph2021} --- 5/8 cumplidos, 2 parciales, 1 no cumplido.}
\label{tab:ralph-general}
\end{table}

\subsection{Anexo E --- Matriz de trazabilidad completa}
\label{anexo:matriz}

La Tabla~\ref{tab:matriz-completa} reproduce íntegramente \texttt{docs/trazabilidad/matriz.csv} (15
filas, 9 columnas), validada por \texttt{scripts/validate-traceability.sh} y por un parseo estricto con
\texttt{csv.reader} de Python que confirma que las 15 filas tienen exactamente 9 campos cada una ---
cinco de ellas (RF-05, RF-07, RF-08, RF-09, RF-10) tenían comas sin escapar en la columna de evidencia
empírica que rompían el conteo de columnas; se corrigieron entre comillas en la corrección del 2026-09-05
sin alterar el contenido de la evidencia. La columna ``Módulo'' se abrevia en algunas filas por espacio;
el CSV original tiene el nombre completo.

\begin{landscape}
\scriptsize
\setlength{\tabcolsep}{3pt}
\begin{longtable}{p{1.0cm}p{1.0cm}p{2.1cm}p{1.1cm}p{2.8cm}p{1.7cm}p{6.2cm}p{2.0cm}}
\caption{Matriz de trazabilidad completa (\texttt{docs/trazabilidad/matriz.csv}).}\label{tab:matriz-completa}\\
\toprule
\textbf{ID Req.} & \textbf{HU} & \textbf{M\'odulo} & \textbf{MoSCoW} & \textbf{Caso de prueba} & \textbf{Estado} & \textbf{Evidencia emp\'irica} & \textbf{Tipo acceso} \\
\midrule
\endfirsthead
\toprule
\textbf{ID Req.} & \textbf{HU} & \textbf{M\'odulo} & \textbf{MoSCoW} & \textbf{Caso de prueba} & \textbf{Estado} & \textbf{Evidencia emp\'irica} & \textbf{Tipo acceso} \\
\midrule
\endhead
\bottomrule
\endfoot
RF-\allowbreak{}01 & HU-\allowbreak{}01 & Seguridad /\allowbreak{} Auth & Must & \seqsplit{JwtTokenProviderTest.java;Phase4FeaturesTest.java} & VERIFICADO & docs/\allowbreak{}mediciones/\allowbreak{}sec/\allowbreak{}owasp/\allowbreak{}OWASP-\allowbreak{}AUDIT.\allowbreak{}md (control A02) & CRUD-\allowbreak{}ORM \\
RF-\allowbreak{}02 & HU-\allowbreak{}02 & Solicitudes & Must & \seqsplit{SubmissionServiceImplTest.java} & VERIFICADO & JaCoCo 65\% cobertura de lineas real (docs/\allowbreak{}mediciones/\allowbreak{}jacoco/\allowbreak{}COVERAGE.\allowbreak{}md) + verificado end-\allowbreak{}to-\allowbreak{}end contra Docker (docs/\allowbreak{}basedatos/\allowbreak{}CATALOGO-\allowbreak{}SP.\allowbreak{}md) & CRUD-\allowbreak{}ORM + SP (sp\_\allowbreak{}generar\_\allowbreak{}codigo\_\allowbreak{}expediente al crear el perfil Estudiante) \\
RF-\allowbreak{}03 & HU-\allowbreak{}03 & Jurados & Must & \seqsplit{JuradoServiceImplTest.java} & VERIFICADO & JaCoCo 53\% cobertura de lineas real (docs/\allowbreak{}mediciones/\allowbreak{}jacoco/\allowbreak{}COVERAGE.\allowbreak{}md) & SP (sp\_\allowbreak{}asignar\_\allowbreak{}jurado\_\allowbreak{}masivo) \\
RF-\allowbreak{}04 & HU-\allowbreak{}04 & Cronograma & Must & \seqsplit{CronogramaServiceImplTest.java} & VERIFICADO & JaCoCo 48\% cobertura de lineas real (docs/\allowbreak{}mediciones/\allowbreak{}jacoco/\allowbreak{}COVERAGE.\allowbreak{}md) + verificado end-\allowbreak{}to-\allowbreak{}end contra Docker & CRUD-\allowbreak{}ORM + SP (sp\_\allowbreak{}validar\_\allowbreak{}conflicto\_\allowbreak{}jurado) \\
RF-\allowbreak{}05 & HU-\allowbreak{}05 & Evaluaciones & Must & \seqsplit{RubricEvaluationServiceImplTest.java} & VERIFICADO & ninguna adicional (cita previa de k6 era incorrecta,\allowbreak{} ver historias/\allowbreak{}HU-\allowbreak{}05.\allowbreak{}md) & CRUD-\allowbreak{}ORM \\
RF-\allowbreak{}06 & HU-\allowbreak{}06 & Actas & Must & \seqsplit{MinutesServiceImplTest.java} & VERIFICADO & Postman Test Collection + PDF real generado con iText en directorio temporal durante el test (verificado 2026-\allowbreak{}08-\allowbreak{}29) & CRUD-\allowbreak{}ORM \\
RF-\allowbreak{}07 & HU-\allowbreak{}07 & Actas /\allowbreak{} Firmas & Should & \seqsplit{MinutesServiceImplTest.java} & VERIFICADO & docs/\allowbreak{}mediciones/\allowbreak{}sec/\allowbreak{}owasp/\allowbreak{}OWASP-\allowbreak{}AUDIT.\allowbreak{}md (control A01,\allowbreak{} verificacion manual de acceso) + docs/\allowbreak{}basedatos/\allowbreak{}CATALOGO-\allowbreak{}SP.\allowbreak{}md & SP (sp\_\allowbreak{}firmar\_\allowbreak{}acta\_\allowbreak{}digital) \\
RF-\allowbreak{}08 & HU-\allowbreak{}08 & Notificaciones & Could & \seqsplit{NotificacionServiceImplTest.java} & VERIFICADO & Postman Test Collection (verificacion manual) + NotificacionServiceImplTest.\allowbreak{}java (11 @Test,\allowbreak{} verificado 2026-\allowbreak{}09-\allowbreak{}05) & CRUD-\allowbreak{}ORM \\
RF-\allowbreak{}09 & HU-\allowbreak{}09 & Reportes & Could & ninguna & NO\_\allowbreak{}VERIFICADO & ninguna (cita previa de k6 era incorrecta,\allowbreak{} ver casos-\allowbreak{}de-\allowbreak{}uso/\allowbreak{}CU-\allowbreak{}09.\allowbreak{}md) & CRUD-\allowbreak{}ORM \\
RF-\allowbreak{}10 & HU-\allowbreak{}10 & Salas & Should & ninguna \seqsplit{(SalaServiceImplTest.java} no existe) & NO\_\allowbreak{}VERIFICADO & ninguna (cita previa de SUS era un non-\allowbreak{}sequitur,\allowbreak{} ver historias/\allowbreak{}HU-\allowbreak{}10.\allowbreak{}md) & CRUD-\allowbreak{}ORM \\
RF-\allowbreak{}11 & HU-\allowbreak{}11 & Usuarios & Must & \seqsplit{AppUserServiceImplTest.java} & VERIFICADO & JaCoCo 69\% cobertura de lineas real + OWASP-\allowbreak{}AUDIT.\allowbreak{}md control A01 & CRUD-\allowbreak{}ORM \\
RF-\allowbreak{}12 & HU-\allowbreak{}12 & Anteproyectos & Must & \seqsplit{AnteproyectoServiceImplTest.java} & VERIFICADO & JaCoCo 72\% cobertura de lineas real + STATIC-\allowbreak{}ANALYSIS.\allowbreak{}md (UNSAFE\_\allowbreak{}HASH\_\allowbreak{}EQUALS corregido) & CRUD-\allowbreak{}ORM \\
RF-\allowbreak{}13 & HU-\allowbreak{}13 & Reportes & Should & \seqsplit{ReporteServiceImplTest.java} & VERIFICADO & Verificado end-\allowbreak{}to-\allowbreak{}end contra Docker (resumen/\allowbreak{}solicitudes-\allowbreak{}por-\allowbreak{}estado/\allowbreak{}sustentaciones-\allowbreak{}por-\allowbreak{}periodo/\allowbreak{}actas/\allowbreak{}actividad-\allowbreak{}docente/\allowbreak{}por-\allowbreak{}carrera devuelven 200 para ADMIN y COORDINADOR;\allowbreak{} 403 para DOCENTE/\allowbreak{}ESTUDIANTE) -\allowbreak{}-\allowbreak{} ver docs/\allowbreak{}actas/\allowbreak{}HISTORIAL-\allowbreak{}Y-\allowbreak{}REPORTES.\allowbreak{}md & CRUD-\allowbreak{}ORM (COUNT/\allowbreak{}GROUP BY en JPQL) \\
RF-\allowbreak{}14 & HU-\allowbreak{}14 & Actas /\allowbreak{} Gestion & Should & \seqsplit{MinutesServiceImplTest.java} & VERIFICADO & E2E contra Docker: DOCENTE ve solo sus actas;\allowbreak{} GET /\allowbreak{}actas/\allowbreak{}\{id\}/\allowbreak{}historial de un acta ajena -\allowbreak{}> 403 (IDOR);\allowbreak{} FINALIZADA-\allowbreak{}>REVISADA -\allowbreak{}> 400 (transicion invalida);\allowbreak{} OBSERVADA/\allowbreak{}ANULADA sin motivo -\allowbreak{}> 400;\allowbreak{} DOCENTE cambiar estado -\allowbreak{}> 403 & CRUD-\allowbreak{}ORM \\
RF-\allowbreak{}15 & HU-\allowbreak{}15 & Actas /\allowbreak{} Historial-\allowbreak{}Auditoria & Should & \seqsplit{MinutesServiceImplTest.java} & VERIFICADO & Migracion V19 crea presus.\allowbreak{}historial\_\allowbreak{}estados\_\allowbreak{}acta (FK a actas/\allowbreak{}estados\_\allowbreak{}acta/\allowbreak{}usuarios) + siembra 10.\allowbreak{}783 eventos CREAR;\allowbreak{} cada cambio de estado inserta una fila con usuario+rol+estado\_\allowbreak{}anterior+estado\_\allowbreak{}nuevo+motivo+fecha;\allowbreak{} la auditoria generica de V15 (trigger trg\_\allowbreak{}auditoria\_\allowbreak{}actas) se conserva como respaldo & CRUD-\allowbreak{}ORM + Trigger (fn\_\allowbreak{}auditoria\_\allowbreak{}generica de V15) \\
\end{longtable}
\end{landscape}
